\chapter{Reproduction}
\label{app:repro}

\section{Environment}

\begin{small}
\begin{verbatim}
git clone https://github.com/Looong01/Splatt3R-SLAM.git
cd Splatt3R-SLAM
conda create -n splatt3r-slam python=3.11 -y
conda activate splatt3r-slam
pip install torch torchvision --index-url \
    https://download.pytorch.org/whl/cu132
pip install cmake ninja swig mkl-devel
pip install -r requirements.txt
pip install -e thirdparty/in3d
pip install --no-build-isolation thirdparty/asmk
\end{verbatim}
\end{small}

\noindent
faiss is built from the vendored source with GPU support; the prebuilt PyPI
wheels must not be installed, because they shadow the build and pin
\texttt{numpy<2}:

\begin{small}
\begin{verbatim}
cd thirdparty/faiss
cmake -B build . && make -C build -j$(nproc) faiss swigfaiss
pip install --no-build-isolation --no-deps ./build/faiss/python
cd ../..
export CUDA_HOME=/usr/local/cuda
export MAX_JOBS=$(nproc)
pip install --no-build-isolation thirdparty/lietorch
pip install --no-build-isolation \
    thirdparty/diff-gaussian-rasterization-modified
pip install --no-build-isolation -e .
cd splatt3r_core/src/mast3r_src/dust3r/croco/models/curope
python setup.py build_ext --inplace
\end{verbatim}
\end{small}

\noindent
The last step builds the CUDA RoPE2D kernel. It is optional --- a PyTorch
fallback exists --- but the fallback is materially slower, and the vendored
\texttt{kernels.cu} in this repository already carries the two changes that
newer PyTorch requires (\texttt{tokens.type()} $\rightarrow$
\texttt{tokens.scalar\_type()}, plus a BFloat16 dispatch), so it builds
without patching.

\section{Running the system}

\begin{small}
\begin{verbatim}
python main.py --dataset datasets/tum/rgbd_dataset_freiburg1_desk \
    --config config/base.yaml --no-viz \
    --refiner --refiner-gpu 1 --refiner-duty 1.0 \
    --refiner-polish-secs 300 --refiner-freeze-means \
    --refiner-conf-fade 0
\end{verbatim}
\end{small}

\noindent
This example selects the released head, fixes centres, disables insertion
thinning and requests a 300\,s polish. Family-head and thinning studies
must specify their checkpoint and attenuation options separately. Exact
saved maps used in the external comparison are identified by hashes in the
evaluation evidence; rerunning this example is not a claim of bit-identical
multi-process reconstruction.

Rendering to disk for figures and the demonstration animations uses
\texttt{--render-dir}. The interactive viewer's \texttt{--save-gs-view} is
written by the visualisation process and therefore produces nothing under
\texttt{--no-viz}; this cost one round of empty output directories before it
was understood.

\section{Scoring a map}

The current external comparison uses
\texttt{scripts/paper/render\_comparisons.py}. It retains all exported SH
bands and uses the corrected MonoGS frame association. The older
\texttt{scripts/eval\_map\_quality.py} is the historical DC-only harness
for component studies and should not regenerate the fresh external table.

\begin{small}
\begin{verbatim}
CUDA_VISIBLE_DEVICES=0 python scripts/paper/render_comparisons.py \
    --scenes office0 office1 office2 office3
CUDA_VISIBLE_DEVICES=1 python scripts/paper/render_comparisons.py \
    --scenes desk office4 room0 room1 room2
\end{verbatim}
\end{small}

The commands require the project environment, datasets and saved maps
referenced by the script. Outputs in \texttt{logs/paper\_ral\_20261007}
include hashes, alignment, frame indices, per-frame metrics and representative
renders. All methods use the same non-keyframe cameras and exposure
normalisation. Trajectories are scored with \texttt{evo\_ape tum -as}. In a headless
environment \texttt{evo} must be told to use a non-interactive backend
(\texttt{evo\_config set plot\_backend Agg}) or it fails at import.

\section{The self-consistency check}

Before a cross-system number is reported, the map is rendered from the poses
used to build it. This check found an inverted camera-to-world convention in
the MonoGS adapter that produced $6.09$\,dB instead of $13.85$\,dB
(\cref{sec:protocol:selfcheck}).

\section{Figures and independently formatted documents}

The figure generators read saved JSON, PNG, and map evidence.
The method overviews are native PowerPoint shapes with separate RA-L and MSc
equation references. The older TikZ figure records the process and device
layout. Captions, prose, numbering, classes, and bibliography styles remain
independent.

\begin{small}
\begin{verbatim}
python scripts/paper/make_ral_figures.py
CUDA_VISIBLE_DEVICES=0 python \
  scripts/paper/render_refinement_pair.py
PYTHONPATH=/path/to/python-pptx python \
  scripts/paper/make_overview_pptx.py
cd docs/Thesis
./build.sh master
./build.sh ral
./build.sh review
\end{verbatim}
\end{small}

The MSc document uses the A4 report class, 12-point text, 1.5 line spacing,
chapter-based structure and numeric \texttt{unsrtnat} bibliography.
The RA-L reading draft uses IEEEtran; the anonymous review draft uses
ieeeconf. All three use pdfLaTeX and distribution-bundled fonts.
The common scientific assets do not import either document's preamble
or prose into the other.
